\documentclass{controle}
\usepackage{main}

\title{Évaluation n°15 : Démonstrations par l'absurde}
\author{Seconde 3}
\date{27 Mars 2026}

% \printanswers

\begin{document}
\titleandname{}
\version
\begin{questions}
\question On donne ci-dessous la démonstration sur l'existence d'une infinité de nombres premiers \textbf{dans le désordre}. Remettre dans l'ordre la démonstration.
\begin{parts}
\part\label{premier4} On a abouti à une \textbf{contradiction}. On en déduit qu'il y a une infinité de nombres premiers.
\part\label{premier2} On pose le nombre $P = p_1 \times p_2 \times p_3 \times \dots \times p_n + 1$. On ne peut pas factoriser $P$ par $p_1$, ou par $p_2$, ou par $p_3$, \dots, ou par $p_n$ (à cause du $+1$). \textbf{On en déduit que $P$ est premier.}
\part\label{premier1} On suppose \textbf{par l'absurde} qu'il y a un nombre fini de nombres premiers. On suppose qu'il y en a $n$. On peut donc établir une liste $p_1; p_2; \dots; p_n$ de nombres premiers.
\part\label{premier3} Or, $P$ est strictement supérieur à $p_1$; $p_2$; $p_3$; \dots; $p_n$. Il ne fait donc clairement pas partie de la liste des nombres premiers $p_1; p_2; \dots; p_n$. \textbf{On en déduit que $P$ n'est pas premier.}
\end{parts}
\begin{solution}[0.8cm]
\begin{itemize}
\item \ref{premier1}
\item \ref{premier2}
\item \ref{premier3}
\item \ref{premier4}
\end{itemize}
\end{solution}
\question On donne ci-dessous la démonstration que $\dfrac{1}{9}$ n'est pas décimal \textbf{dans le désordre}. Remettre dans l'ordre la démonstration.
\begin{parts}
\part\label{deci2} On simplifie l'égalité et on obtient $10^k = 9a$. \textbf{Donc $10^k$ est un multiple de $9$}. 
\part\label{deci3} Or, le nombre $10^k$ s'écrit avec un seul $1$ et $k$ de $0$. Donc la somme de ses chiffres fait $1$. Cette somme n'est pas un multiple de $9$. \textbf{Donc $10^k$ n'est pas un multiple de $9$}.
\part\label{deci1} On suppose \textbf{par l'absurde} que $\dfrac{1}{9}$ est décimal. Il existe donc $a \in \Z$ et $k \in \N$ tel que $\dfrac{1}{9}=\dfrac{a}{10^k}$.
\part\label{deci4} On a abouti à une \textbf{contradiction}. On en déduit que $\dfrac{1}{9}$ n'est pas décimal.
\end{parts}
\begin{solution}[0.8cm]
\begin{itemize}
\item \ref{deci1} 
\item \ref{deci2} 
\item \ref{deci3} 
\item \ref{deci4} 
\end{itemize}
\end{solution}
\question On donne ci-dessous la démonstration que $\sqrt{3}$ est irrationnel \textbf{dans le désordre}. Remettre dans l'ordre la démonstration.
\begin{parts}
\part\label{sqrt5} On a abouti à une \textbf{contradiction}. On en déduit que $\sqrt{3}$ est irrationnel.
\part\label{sqrt3} Puisque $p$ est un multiple de $3$, il est de la forme $3k$. On obtient donc $(3k)^2 = 3q^2$ qui se simplifie en $3k^2 = q^2$. Par le même raisonnement que précédemment, on obtient que \textbf{$q$ est un multiple de $3$}.
\part\label{sqrt2} On simplifie l'égalité pour obtenir $p^2 = 3q^2$. Or, on utilise un résultat admis que si un nombre $n^2$ est un multiple de $3$, alors $n$ est un multiple de $3$. \textbf{Donc $p$ est un multiple de $3$}. 
\part\label{sqrt4} Ainsi, $p$ et $q$ étant des multiples de $3$, \textbf{la fraction $\dfrac{p}{q}$ n'est pas irréductible}.
\part\label{sqrt1} On suppose par l'absurde que $\sqrt{3}$ est rationnel. Donc il existe une fraction irréductible $\dfrac{p}{q}$ telle que $\sqrt{3} = \dfrac{p}{q}$.
\end{parts}
\begin{solution}[0.8cm]
\begin{itemize}
\item \ref{sqrt1} 
\item \ref{sqrt2} 
\item \ref{sqrt3} 
\item \ref{sqrt4} 
\item \ref{sqrt5} 
\end{itemize}
\end{solution}
\end{questions}





\newpage
\titleandname{}
\version





\begin{questions}
\question On donne ci-dessous la démonstration que $\dfrac{1}{9}$ n'est pas décimal \textbf{dans le désordre}. Remettre dans l'ordre la démonstration.
\begin{parts}
\part\label{deci3} Or, le nombre $10^k$ s'écrit avec un seul $1$ et $k$ de $0$. Donc la somme de ses chiffres fait $1$. Cette somme n'est pas un multiple de $9$. \textbf{Donc $10^k$ n'est pas un multiple de $9$}.
\part\label{deci4} On a abouti à une \textbf{contradiction}. On en déduit que $\dfrac{1}{9}$ n'est pas décimal.
\part\label{deci1} On suppose \textbf{par l'absurde} que $\dfrac{1}{9}$ est décimal. Il existe donc $a \in \Z$ et $k \in \N$ tel que $\dfrac{1}{9}=\dfrac{a}{10^k}$.
\part\label{deci2} On simplifie l'égalité et on obtient $10^k = 9a$. \textbf{Donc $10^k$ est un multiple de $9$}. 
\end{parts}
\begin{solution}[0.8cm]
\begin{itemize}
\item \ref{deci1} 
\item \ref{deci2} 
\item \ref{deci3} 
\item \ref{deci4} 
\end{itemize}
\end{solution}
\question On donne ci-dessous la démonstration sur l'existence d'une infinité de nombres premiers \textbf{dans le désordre}. Remettre dans l'ordre la démonstration.
\begin{parts}
\part\label{premier2} On pose le nombre $P = p_1 \times p_2 \times p_3 \times \dots \times p_n + 1$. On ne peut pas factoriser $P$ par $p_1$, ou par $p_2$, ou par $p_3$, \dots, ou par $p_n$ (à cause du $+1$). \textbf{On en déduit que $P$ est premier.}
\part\label{premier3} Or, $P$ est strictement supérieur à $p_1$; $p_2$; $p_3$; \dots; $p_n$. Il ne fait donc clairement pas partie de la liste des nombres premiers $p_1; p_2; \dots; p_n$. \textbf{On en déduit que $P$ n'est pas premier.}
\part\label{premier1} On suppose \textbf{par l'absurde} qu'il y a un nombre fini de nombres premiers. On suppose qu'il y en a $n$. On peut donc établir une liste $p_1; p_2; \dots; p_n$ de nombres premiers.
\part\label{premier4} On a abouti à une \textbf{contradiction}. On en déduit qu'il y a une infinité de nombres premiers.
\end{parts}
\begin{solution}[0.8cm]
\begin{itemize}
\item \ref{premier1}
\item \ref{premier2}
\item \ref{premier3}
\item \ref{premier4}
\end{itemize}
\end{solution}
\question On donne ci-dessous la démonstration que $\sqrt{3}$ est irrationnel \textbf{dans le désordre}. Remettre dans l'ordre la démonstration.
\begin{parts}
\part\label{sqrt3} Puisque $p$ est un multiple de $3$, il est de la forme $3k$. On obtient donc $(3k)^2 = 3q^2$ qui se simplifie en $3k^2 = q^2$. Par le même raisonnement que précédemment, on obtient que \textbf{$q$ est un multiple de $3$}.
\part\label{sqrt5} On a abouti à une \textbf{contradiction}. On en déduit que $\sqrt{3}$ est irrationnel.
\part\label{sqrt2} On simplifie l'égalité pour obtenir $p^2 = 3q^2$. Or, on utilise un résultat admis que si un nombre $n^2$ est un multiple de $3$, alors $n$ est un multiple de $3$. \textbf{Donc $p$ est un multiple de $3$}. 
\part\label{sqrt1} On suppose par l'absurde que $\sqrt{3}$ est rationnel. Donc il existe une fraction irréductible $\dfrac{p}{q}$ telle que $\sqrt{3} = \dfrac{p}{q}$.
\part\label{sqrt4} Ainsi, $p$ et $q$ étant des multiples de $3$, \textbf{la fraction $\dfrac{p}{q}$ n'est pas irréductible}.
\end{parts}
\begin{solution}[0.8cm]
\begin{itemize}
\item \ref{sqrt1} 
\item \ref{sqrt2} 
\item \ref{sqrt3} 
\item \ref{sqrt4} 
\item \ref{sqrt5} 
\end{itemize}
\end{solution}
\end{questions}
\end{document}